\chapter{Metodologia de Pesquisa}

A metodologia explica o caminho utilizado para produzir conhecimento. Ela deve ser compatível com o problema, os objetivos e a modalidade do TCC. Um método bem descrito permite que o leitor compreenda como as informações foram obtidas, organizadas e analisadas.

\section{Abordagens de pesquisa}

A abordagem qualitativa busca compreender sentidos, experiências, processos, relações, discursos e contextos. A abordagem quantitativa trabalha com mensuração, frequência, associação, comparação estatística ou descrição numérica. A abordagem mista combina procedimentos qualitativos e quantitativos.

A escolha da abordagem não deve ser feita por preferência pessoal, mas pela natureza do problema. Perguntas sobre percepções, significados e processos costumam exigir abordagem qualitativa. Perguntas sobre prevalência, frequência e associação podem exigir abordagem quantitativa. Problemas complexos podem se beneficiar de desenho misto.

\section{Pesquisa bibliográfica}

A pesquisa bibliográfica utiliza produção acadêmica já publicada como base principal. Deve indicar bases, descritores, critérios de seleção, período, tipos de documento e estratégia de análise.

Uma pesquisa bibliográfica frágil apenas reúne textos encontrados. Uma pesquisa bibliográfica consistente seleciona fontes de modo justificado, compara argumentos e constrói síntese crítica.

\section{Pesquisa documental}

A pesquisa documental analisa documentos que podem incluir legislação, relatórios, atas, políticas públicas, registros institucionais, processos, planos, diretrizes, imagens, materiais de divulgação, dados administrativos ou arquivos históricos.

O estudante deve explicar origem dos documentos, critérios de seleção, confiabilidade, contexto de produção e forma de análise.

\section{Pesquisa de campo}

A pesquisa de campo envolve aproximação direta com pessoas, grupos, instituições ou ambientes. Pode utilizar entrevistas, questionários, observação, grupos focais, registros fotográficos, diário de campo ou instrumentos específicos da área.

Esse tipo de pesquisa exige planejamento cuidadoso, autorização quando necessária, atenção ética e registro organizado.

\section{População, amostra e corpus}

O trabalho deve identificar de onde provêm as evidências analisadas. Em pesquisas com pessoas, isso inclui população de interesse, critérios de inclusão e exclusão, forma de recrutamento, tamanho e composição da amostra. Em pesquisas documentais ou bibliográficas, inclui a delimitação do corpus, das bases, do período e dos tipos de documento.

Quando houver intenção de generalizar resultados, a estratégia amostral precisa sustentar essa pretensão. Amostras por conveniência podem ser úteis, mas suas limitações devem ser reconhecidas. Distribuição etária, territorial ou de outros grupos relevantes deve ser comparada ao público sobre o qual se pretende concluir.

\section{Entrevistas}

Entrevistas podem ser estruturadas, semiestruturadas ou abertas. O roteiro deve estar alinhado aos objetivos e evitar perguntas que induzam respostas.

Na análise, o estudante não deve apenas transcrever falas. É necessário interpretar os dados, organizar categorias, relacionar respostas ao problema e dialogar com a literatura.

\section{Questionários}

Questionários podem coletar dados de modo padronizado. Devem usar perguntas claras, linguagem adequada ao público, alternativas coerentes e ordem lógica. Antes da aplicação, recomenda-se testar o instrumento com pequeno grupo.

Dados de questionários precisam ser apresentados e analisados com cuidado. Percentuais sem contexto podem gerar interpretações equivocadas, especialmente em amostras pequenas.

\section{Validade e controle de qualidade da coleta}

O planejamento deve prever teste do instrumento, treinamento da equipe quando necessário, registro de recusas e perdas, prevenção de duplicidade e critérios para tratar respostas incompletas, inconsistentes ou suspeitas. Em formulários digitais, controles contra automação e múltiplos envios podem ser necessários, desde que não coletem identificadores excessivos.

Autenticidade não significa necessariamente conhecer a identidade civil do participante. O objetivo é reunir evidências suficientes de elegibilidade e unicidade de acordo com o risco do estudo, preservando a privacidade. Procedimentos de validação devem ser descritos na metodologia e compatíveis com o protocolo ético.

\section{Observação}

A observação pode ser participante ou não participante, estruturada ou aberta. Deve indicar local, período, foco de observação, forma de registro e cuidados éticos.

O diário de campo é recurso importante para registrar situações, impressões, eventos e reflexões. Esses registros devem ser tratados como material de análise, não como memória informal.

\section{Análise de dados}

A análise transforma material bruto em resultados interpretáveis. Pode envolver categorização, análise temática, análise de conteúdo, estatística descritiva, comparação documental, análise de discurso, triangulação ou outro procedimento adequado.

O estudante deve explicar como passou dos dados aos resultados. A banca precisa enxergar o caminho analítico.

\section{Plano de gestão dos dados}

Pesquisas empíricas devem indicar como os dados serão nomeados, armazenados, protegidos, versionados e preparados para análise. Quando houver dados pessoais, o plano também deve definir perfis de acesso, separação de identificadores, prazo de retenção e forma de descarte ou compartilhamento.

O conjunto bruto não deve ser alterado silenciosamente. Exclusões, correções, transformações e critérios de tratamento precisam ser registrados para que o percurso analítico possa ser auditado.

\section{Limitações metodológicas}

Todo trabalho possui limites. Reconhecê-los não enfraquece o TCC; ao contrário, demonstra maturidade. Limitações podem envolver amostra, acesso a dados, tempo, escopo, instrumentos, generalização ou disponibilidade de fontes.

\caixaatencao{Método não é decoração textual. Se a metodologia não explica como o trabalho foi realizado, o leitor não consegue avaliar a validade dos resultados.}
